\begingroup
\makeatletter\tagsleft@true\let\veqno\@@leqno\makeatother
\fancyhead{}\fancyhead[R]{\small279}\fancyhead[C]{\scriptsize BEMERKUNGEN ZUM AGM. ABSCHNITT VIII.}
\SourceLeaf{286b}
\pdfbookmark[1]{Erläuterungen zu [VIII.]}{x1-agm-erlaeuterung-viii}
\begin{center}\textls[120]{Erläuterungen zu [VIII.], S. 213—216.}\end{center}
\par\indent Die durch [1] eingeführte Größe $x$ ist mit der in der Scheda Ac (siehe Abschnitt [V.], Gleichung [18], S. 197) durch $z$ bezeichneten Größe identisch; es ist
\[x^4=e^{-\pi\frac{M(a,b)}{M(a,c)}}.\]
\noindent Die Darstellung [1] ergibt sich durch Grenzübergang bei wiederholter Anwendung der \textsc{Landen}schen Transformation
\[\sin\psi=\operatorname{tg}^2\frac\varphi2\]
\noindent vergl. oben S. 198), sie findet sich auch in der Abhandlung aus Handbuch 18, Bd, Werke III, S. 448, wo aber die hier mit $x$ bezeichnete Größe mit $\sqrt[4]x$ bezeichnet ist. Die Formeln [2], wo $h=M(a,b)$ ist, sind uns schon von den Schedae Ac und Af her bekannt. Das wesentlich Neue, was uns hier entgegentritt, ist der Algorithmus [3]. \textsc{Gauss} erhält ihn anscheinend durch folgenden Gedankengang.
\par\indent Die Gleichungen [5] lauten in neuerer Bezeichnungsweise (vergl. oben S. 275)
\begin{equation*}\tag*{$[5]'$}\begin{gathered}
A=H\,.\,\vartheta_{00}^2(v|q,\quad B=H\,.\,\vartheta_{01}^2\,v|q);\quad v=e^{2\pi iy},\ q=x^4.
\end{gathered}\end{equation*}
\noindent Es ist also
\begin{equation*}\tag*{$[5]''$}\frac BA=\frac{\vartheta_{01}^2(v|q)}{\vartheta_{01}^2(v|q)}.\end{equation*}
\noindent Ähnlich, wie aus
\[\frac ba=\frac{\vartheta_{01}^2(0|q)}{\vartheta_{00}^2(0|q)}\]
\noindent gebildet ist
\[\frac{b'}{a'}=\frac{\vartheta_{01}^2(0|q^2)}{\vartheta_{00}^2(0|q^2)},\]
\noindent bildet nun \textsc{Gauss} aus $[5]''$
\[\frac{B'}{A'}=\frac{\vartheta_{01}^2(2v|q^2)}{\vartheta_{00}^2(2v|q^2)}.\]
\noindent Nach den in der Scheda Ac enthaltenen Formeln für die Transformation zweiter Ordnung (siehe oben die Gleichungen $[14]'$, $[17]'$, S. 275)
\clearpage
\ArithmeticHead{280}{L}{ANALYSIS. NACHLASS.}
\SourceLeaf{287a}
\[\begin{aligned}
2\vartheta_{00}(0|q^2)\,\vartheta_{00}(2v|q^2)&=\vartheta_{00}^2(v\,q)+\vartheta_{01}^2(v|q),\\
\vartheta_{01}(0|q^2)\,\vartheta_{01}(2v|q^2)&=\vartheta_{00}(v|q)\,\vartheta_{01}(v|q)
\end{aligned}\]
\noindent ergibt sich aus den obigen Formeln
\begin{equation*}\tag*{$[5]'''$}\frac{B'}{A'}=\frac{AB}{\left(\frac{A+B}2\right)^2}\frac{a'}{b'}.\end{equation*}
\noindent \textsc{Gauss} setzt nun fest, daß wie beim agM.
\[A'=\frac{A+B}2\]
\noindent sein soll: dann folgt aber aus $[5]'''$
\[B'=\frac{2AB}{A+B}\frac{a'}{b'}.\]
\par\indent Auf eine weitere Erläuterung der \textsc{Gauss}schen Formeln können wir verzichten, da sie mit Hilfe der Theorie der Thetafunktionen leicht bestätigt werden können, wir bemerken nur noch, daß die in [6] festgesetzte Beziehung zwischen den aufeinanderfolgenden Werten $\varphi,\varphi',\ldots$ sich in der Form
\begin{equation*}\tag*{$[6]'$}
\sin\varphi=\frac ac\frac{2\sin\varphi'\sqrt{\frac{c'}{a'}}}{1+\sin^2\varphi'\,.\,\frac{c'}{a'}}=\frac{2a\sin\varphi'}{(a+b)+(a-b)\sin^2\varphi'}
\end{equation*}
\noindent darstellen läßt\textsuperscript{*)}. Es ist dies die sogenannte \textsc{Gauss}sche Transformation der \textit{Determinatio attractionis} (1818), Werke III, S. 352. Die Darstellung von $\sin\varphi$ und $\operatorname{tang}\varphi$ durch Thetafunktionen in den Gleichungen [8] und [9] stimmt — von geringfügigen Schreibfehlern abgesehen, die beim Abdruck oben S. 216 beseitigt sind — mit der aus dem Jahre 1827 stammenden im Handbuch 16, Bb (abgedruckt Werke III, S. 473, art. [5.]) überein, wenn man das dortige $\vartheta=\varphi-\frac\pi2$ setzt. Bemerkenswert ist, daß oben S. 216, Gleichung [9] (im Nenner von $\operatorname{tang}\varphi$) auch die ungerade Thetafunktion in derselben Form wie a. a. O. S. 472 auftritt\textsuperscript{**)}.
\par\indent Die Aufzeichnung [VIII.] beginnt auf S. 37 der Scheda An; auf S. 35 dieser Scheda findet sich am Schluß einer astronomischen Rechnung die Bemerkung: »geendigt d. 2. May 1809«. Unser Abschnitt [VIII.] ist danach mit der Tagebuchnotiz Nr. 139 vom 20. Juni 1809 »Series ad Media arithmetico-geometrica pertinentia fusius evolutae« in Verbindung zu bringen. Vergl. auch die aus dem Handbuch 18, Bd, S. 221—233 stammende Abhandlung, die Werke III, S. 446 abgedruckt ist, und (vergl. die Bemerkung \textsc{Scherings} ebenda S. 494) aus derselben Zeit stammt.
\Needspace{4\baselineskip}
\par\smallskip\noindent\rule{.25\linewidth}{.4pt}\par
{\small
\par\indent *) Es ist nämlich $\sin2M=\frac ca$; vergl. auch die aus Scheda Ac herrührende Formel $[14]'$, S. 275.
\par\indent **) \textsc{Schering} bemerkt Werke III, S. 390 (wohl in Bezug auf die vorstehende Aufzeichnung), daß die Bestimmung eines reellen $y$ »durch eine Lücke unvollendet gelassen« sei. In der Tat befindet sich (siehe oben S. 214, nach Gl. [4]) an der Stelle, wo die Erklärung von $y$ folgen sollte, eine Lücke, die jedoch durch einen Hinweis auf die im art. [2.] folgende Gleichung [7] ausgefüllt werden konnte. Vergl. auch die im Abschnitt [IX.] artt. [10.—12.], S. 227 ff. wiedergegebenen Aufzeichnungen.
\par}
\clearpage
\endgroup
